\section{Related Works}
\textbf{Optimization methods.} Many previous efforts have focused on optimization objectives that target classifier performance metrics to try to balance precision and recall in some fashion~\citep{rakotomamonjyOptimizingAreaRoc2004,burgesLearningRankNonsmooth2006,yueSupportVectorMethod2007,liptonOptimalThresholdingClassifiers2014}.
As we review below, most of these methods either pursue different objectives (less appropriate for our use case than Eq.~\eqref{eq:max_recall_st_precision}), have limited scalability to large datasets, or have deficiencies in approximation quality.

Optimizing area under the precision-recall curve (AURPC) has been pursued by  
\citet{metzlerMarkovRandomField2005} and \citet{caruanaEmpiricalComparisonSupervised2006}, as well as recently by \citet{ramzi2021robust} and \citet{qi2021stochastic}, who both propose tractable surrogate losses for gradient-based optimization.
However, optimizing AUPRC cannot maximize the quality of recall at a specific precision value, as our later method can.
Figure \ref{fig:roc_and_pr_curves_eICU} provides a real-data example of where one method dominates in AUPRC, but another (ours) achieves a better recall at the desired false alarm rate.
%showing the PR curves on the validation and test sets of the e-ICU dataset. In these plots, BCE has a better AUPRC, but our sigmoid bound archives a better recall at the desired false alarm rate of 0.2, as indicated by the chosen threshold.}
Similar concerns exist for approaches that optimize the F-measure, such as the sigmoid approximations of~\citet{tsoiHeavisideFunctionApproximation2021} or the pseudo-linear methods of \citet{narasimhan2015optimizing} and \citet{puthiya2014optimizing}.
% none of them pursue our intended objective (Eq.~\eqref{eq:max_recall_st_precision}). 

\citet{fathonyAPperfIncorporatingGeneric2020}'s adversarial prediction (AP) framework allows gradient-based training for a variety of non-decomposable objectives (such as our Eq.~\eqref{eq:max_recall_st_precision}).
However, their AP method suffers from scalability issues, with runtime \emph{cubic} in the number of examples in each minibatch.
Our work is conceptually simpler and keeps runtime \emph{linear} in the number of examples. 
Scalability is also a concern for the methods of \citet{joachimsSupportVectorMethod2005}, which pursue optimizing precision at fixed recall at cost quadratic in the number of examples.

We build upon the work of \citet{ebanScalableLearningNonDecomposable2017}, which developed a framework for gradient-based learning of non-decomposable losses like our objective in Eq.~\eqref{eq:max_recall_st_precision}.
\citeauthor{ebanScalableLearningNonDecomposable2017} accomplish this via tractable hinge-loss bounds of the counts of true and false positives (as detailed in Sec.~\ref{sec:methods_hinge}).
However, we find in practice these hinge bounds are so loose (see Fig.~\ref{fig:tp_fp_bounds}) that they cannot solve the problem: solutions often do not satisfy the desired precision constraint (see Fig.~\ref{fig:toy_data_precision_90}).  We will show that our tighter bounds and more careful treatment of the optimization problem lead to noticeably better solutions.
Our surrogate bound ideas originally appeared in a preliminary workshop paper~\citep{rathOptimizingClinicalEarly2021}.
This present paper offers more rigorous conceptual justification as well as substantially expanded empirical evaluation (comparing to more alternative methods as well as more realistic ``continual monitoring'' tasks).


%We develop tighter bounds and show that they can solve real problems with runtime \emph{linear} in the number of examples.

\textbf{Clinical methods.}
Some false alarm control methods have been directly developed for use in a hospital (\citet{chambrin2001alarms}, \citet{antink2016reducing}, \citet{eerikainen2016reduction}).
\citet{au2019reduction} reduce false alarms in the ICU via post-hoc feature selection, but do not encode limits on false alarms into any training objective.
\citet{hever2019machine} reduce false alarms for arrhythmia by training random forests to match expert rules. However, such rules are not easily available for all tasks.
Our work allows direct control of the desired precision level and applies to many models and tasks.


%\citet{au2019reduction} suggest a method to reduce false  alarms in the ICU using EKG, PPG and BP signals by extracting signal processing based sensor-specific features, and ranking them by their importance using a random forest, before classifying arrhythmias. They report competitive performance in terms of a "score", which is proportional to false alarm rates, but has no direct relation to it. Moreover, their features extraction pipeline is specialized for high frequency physiological sensor data, which might not always be available. 

%\citet{hever2019machine} demonstrate reduction in false arrhythmia alarms in the ICU by training a random forest in a data-driven strategy using clinical expert based rules. They built a bank of alarm classifiers for multiple scenarios of missing sensor data, using the full expert rules as ground truth, and found considerable improvement on the Youden's index compared to a set of partial rules. However, they rely on a set of expert based rules, designed specifically for a few physiological sensors and 7 arrhythmia conditions. Obtaining such rules on vast EHR datasets with heterogeneous data isn't feasible. Our work allows the clinician to directly control the desired false alarm rate and can be applied to large EHR datasets with multiple physiological sensors, lab measurements and demographic features.
